\documentclass[a4paper, 12pt]{article}
\usepackage[hmargin=2cm,vmargin=2.5cm]{geometry}

\usepackage{enumerate}
\usepackage{makecell}
\usepackage{graphicx}
\usepackage{hyperref}
\usepackage{amsmath}
\usepackage[backend=biber,style=apa, url=true, sortcites]{biblatex}
\usepackage[table]{xcolor}
\usepackage{minted}
\usepackage{graphicx}


\addbibresource{references.bib}
\hypersetup{
    colorlinks,
    citecolor=black,
    filecolor=black,
    linkcolor=black,
    urlcolor=black
}

\setlength{\arrayrulewidth}{0.4mm}

\newcommand{\HRule}{\rule{\linewidth}{0.5mm}}

\begin{document}

\begin{titlepage}
  \begin{center}
    % logo
    \includegraphics[width=0.5\textwidth]{figures/uhmlogo.png}~\\[1cm]

    \textsc{\Large ATMO XXX \\ Class Name Here \\ Fall 2023}\\[3cm]


    \HRule \\[0.4cm]
    {\large \bfseries Lab 00 \\ Title of the lab \\[0.4cm]}
    \HRule \\[4cm]

    \large\textbf{Student:}\\
    NAME HERE \\ Department of Atmospheric Sciences \\ University of Hawaii at Manoa\\[3cm]

    \vfill

    {\large \today}

  \end{center}
\end{titlepage}

\newpage

\tableofcontents

\newpage

\section{Questions}

\begin{enumerate}[\bfseries Q1]
\item \textbf{Question 1 goes here in bold}
\addcontentsline{toc}{subsection}{Q1}\\\\
answer goes here with a reference as an example (\cite{noauthor_anatomy_nodate}).\\
\begin{enumerate}[a)]
  \item \textbf{Subquestion 1}\\
        Answer goes here
        \\

  \item \textbf{Subquestion 2}\\
        Answer goes here.
        \\

  \item \textbf{Subquestion 3}\\
        Answer goes here

\end{enumerate}

\item \textbf{Question 2 goes here in bold}
\addcontentsline{toc}{subsection}{Q2}\\\\
You can insert code files also

\begingroup
\obeylines
\inputminted[fontsize=\scriptsize, linenos, breaklines=true, xleftmargin=0.75cm, frame=lines]{matlab}{code/MATLAB/test.m}
\endgroup

The first dimension is the number of vertical bins (125) and the second dimension is the number of profiles (37082).
The maximum value for each vertical bin is plot in Figure \ref{fig:max_reflectivity}.

\begin{figure}[htp!]
  \centering
  \includegraphics[width=0.5\textwidth]{figures/reflectivity_max.png}
  \caption{Maximum reflectivity for each vertical bin. The dashed line on bin 105 (0km) correspond to the closest height to mean sea level.}
  \label{fig:max_reflectivity}
\end{figure}

\newpage % insert following content in new page

\item \textbf{Question 3 goes here in bold}
\addcontentsline{toc}{subsection}{Q3}\\\\
answer goes here.\\

\item \textbf{Question 4 goes here in bold}
\addcontentsline{toc}{subsection}{Q4}\\\\
answer goes here.\\

\end{enumerate}

\newpage

\printbibliography

\newpage

\section{Figures}

\begin{figure}[htp!]
  \centering
  \includegraphics[width=0.5\textwidth]{figures/fig1.png}
  \caption{A nice capybara.}
  \label{fig:fig1}
\end{figure}

\begin{figure}[htp!]
  \centering
  \includegraphics[width=0.5\textwidth]{figures/fig1.png}
  \caption{Same capybara}
  \label{fig:fig2}
\end{figure}

\newpage

\section{Code}

\subsection*{Python}
\addcontentsline{toc}{subsection}{Python}

\subsubsection*{test.py}
This code is used to generate figures \ref{fig:fig1} and \ref{fig:fig2}.
\begingroup
\obeylines
\inputminted[fontsize=\scriptsize, linenos, breaklines=true, xleftmargin=0.75cm, frame=lines]{python}{code/Python/test.py}
\endgroup


\end{document}